\documentclass[10pt,a4paper]{amsart}
\usepackage{amsmath,amssymb,amsthm,graphicx}
% Retained preamble for the Stacks Project chapter snapshot in this directory.
% The chapter's own preamble.tex is a copy of the shared Stacks Project preamble; the block
% below is sliced MECHANICALLY, line for line, from it. Nothing was added, removed or
% redefined; the wrapper class, the xr-hyper/hyperref wiring and the externaldocument list
% are omitted because the assembler supplies the document class.
%
% lines taken from preamble.tex: 3, 5, 9, 11, 12, 156, 157, 158, ... (51 lines total)
%
% package coverage restored from the chapter's own preamble (2026-10-13): xy, xy-2cell

\usepackage[all]{xy}
\xyoption{2cell}
\UseAllTwocells

\newenvironment{reference}{\comment}{\endcomment}
\newenvironment{slogan}{\comment}{\endcomment}
\newenvironment{history}{\comment}{\endcomment}
\theoremstyle{plain}
\newtheorem{theorem}[subsection]{Theorem}
\newtheorem{proposition}[subsection]{Proposition}
\newtheorem{lemma}[subsection]{Lemma}
\theoremstyle{definition}
\newtheorem{definition}[subsection]{Definition}
\newtheorem{example}[subsection]{Example}
\newtheorem{exercise}[subsection]{Exercise}
\newtheorem{situation}[subsection]{Situation}
\theoremstyle{remark}
\newtheorem{remark}[subsection]{Remark}
\newtheorem{remarks}[subsection]{Remarks}
\numberwithin{equation}{subsection}
\def\lim{\mathop{\mathrm{lim}}\nolimits}
\def\colim{\mathop{\mathrm{colim}}\nolimits}
\def\Spec{\mathop{\mathrm{Spec}}}
\def\Hom{\mathop{\mathrm{Hom}}\nolimits}
\def\Ext{\mathop{\mathrm{Ext}}\nolimits}
\def\SheafHom{\mathop{\mathcal{H}\!\mathit{om}}\nolimits}
\def\SheafExt{\mathop{\mathcal{E}\!\mathit{xt}}\nolimits}
\def\Sch{\mathit{Sch}}
\def\Mor{\mathop{\mathrm{Mor}}\nolimits}
\def\Ob{\mathop{\mathrm{Ob}}\nolimits}
\def\Sh{\mathop{\mathit{Sh}}\nolimits}
\def\NL{\mathop{N\!L}\nolimits}
\def\CH{\mathop{\mathrm{CH}}\nolimits}
\def\proetale{{pro\text{-}\acute{e}tale}}
\def\etale{{\acute{e}tale}}
\def\QCoh{\mathit{QCoh}}
\def\Ker{\mathop{\mathrm{Ker}}}
\def\Im{\mathop{\mathrm{Im}}}
\def\Coker{\mathop{\mathrm{Coker}}}
\def\Coim{\mathop{\mathrm{Coim}}}
\DeclareMathSymbol{\boxtimes}{\mathbin}{AMSa}{"02}
\def\QCohstack{\mathcal{QC}\!\mathit{oh}}
\def\Cohstack{\mathcal{C}\!\mathit{oh}}
\def\Spacesstack{\mathcal{S}\!\mathit{paces}}
\def\Quotfunctor{\mathrm{Quot}}
\def\Hilbfunctor{\mathrm{Hilb}}
\def\Curvesstack{\mathcal{C}\!\mathit{urves}}
\def\Polarizedstack{\mathcal{P}\!\mathit{olarized}}
\def\Complexesstack{\mathcal{C}\!\mathit{omplexes}}
\def\Pic{\mathop{\mathrm{Pic}}\nolimits}
\def\Picardstack{\mathcal{P}\!\mathit{ic}}
\def\Picardfunctor{\mathrm{Pic}}
\def\Deformationcategory{\mathcal{D}\!\mathit{ef}}
\usepackage{cleveref}
\crefname{originallemma}{original lemma}{original lemmas}
\Crefname{originallemma}{Original lemma}{Original lemmas}
\makeatletter
\newlabel{descent-proposition-equivalence-quasi-coherent}{{03DX}{source}{}{sourcepointer-1}{}}
\newlabel{descent-proposition-equivalence-quasi-coherent@cref}{{[originallemma][0][]03DX}{[1][0][]source}{}{}{}}
\makeatother
\crefname{originallemma}{original lemma}{original lemmas}
\Crefname{originallemma}{Original lemma}{Original lemmas}
\makeatletter
\newlabel{equation-comparison-modules}{{equation-comparison-modules}{source}{}{sourcepointer-2}{}}
\newlabel{equation-comparison-modules@cref}{{[originallemma][0][]equation-comparison-modules}{[1][0][]source}{}{}{}}
\makeatother
\crefname{originallemma}{original lemma}{original lemmas}
\Crefname{originallemma}{Original lemma}{Original lemmas}
\makeatletter
\newlabel{lemma-characterize-quasi-coherent}{{lemma-characterize-quasi-coherent}{source}{}{sourcepointer-3}{}}
\newlabel{lemma-characterize-quasi-coherent@cref}{{[originallemma][0][]lemma-characterize-quasi-coherent}{[1][0][]source}{}{}{}}
\makeatother
\title{Stacks, Sheaves on Algebraic Stacks: lemma-quasi-coherent}
\author{The Stacks Project Authors}
\date{stacks-project:stacks-sheaves.tex:local-snapshot-20261012; GNU FDL 1.2 or later}
\begin{document}
\maketitle
\noindent\textit{Editorial note.} Statement span is lines 1909-1923 of sources/stacks-algebraic-stacks-sheaves/pinned-original-source.tex; proof span is lines 1925-1956, the complete author proof (30 body lines).\par\medskip
\noindent\textit{Editorial note.} Navigation stubs: 3 cited-but-not-included labels are supplied as non-load-bearing pointers with their pinned tags where the snapshot tag index has them.\par\medskip
\noindent\textit{Editorial note.} extra\_packages is empty on purpose: the retained preamble already loads xy with options and a second declaration clashes in the assembler wrapper.\par\medskip
\noindent\textit{Editorial note.} Text is retained verbatim, including the author's own typos and placeholders.\par\medskip
\noindent\textit{Editorial note.} Queue index 15 of queue\_algebraic-stacks-sheaves.json (pb=1687, stmt 1909-1923, proof 1925-1956).\par\medskip
\noindent\textit{Editorial note.} Cross-chapter same-label warning (NOT an occupancy conflict, Chief ruling 2026-10-13 evening): label 'lemma-quasi-coherent' also exists in stacks-schemes-cohomology — a different proposition in a different chapter file.\par\medskip
\noindent\textit{Allowed standard prerequisites (editorial).}\begin{itemize}
\item Explicit prerequisite boundary. The cited inputs are prerequisites; the argument in the retained span is not replaced by any cited result.
\end{itemize}\medskip
\ifdefined\hypertarget\hypertarget{sourcepointer-1}{}\fi
\noindent\textit{Source reference.} Original lemma 03DX: sources/stacks-algebraic-stacks-sheaves/pinned-original-source.tex, source lines 1956--1956. Cited by the retained statement or proof but not part of the collected spans; supplied only as a navigation pointer. The argument does not depend on its content.\par\medskip
\ifdefined\hypertarget\hypertarget{sourcepointer-2}{}\fi
\noindent\textit{Source reference.} Original lemma equation-comparison-modules: sources/stacks-algebraic-stacks-sheaves/pinned-original-source.tex, source lines 1382--1386. Cited by the retained statement or proof but not part of the collected spans; supplied only as a navigation pointer. The argument does not depend on its content.\par\medskip
\ifdefined\hypertarget\hypertarget{sourcepointer-3}{}\fi
\noindent\textit{Source reference.} Original lemma lemma-characterize-quasi-coherent: sources/stacks-algebraic-stacks-sheaves/pinned-original-source.tex, source lines 1798--1805. Cited by the retained statement or proof but not part of the collected spans; supplied only as a navigation pointer. The argument does not depend on its content.\par\medskip
\begin{lemma}
\label{lemma-quasi-coherent}
Let $p : \mathcal{X} \to (\Sch/S)_{fppf}$ be a category
fibred in groupoids. Let $\mathcal{F}$
be a presheaf of $\mathcal{O}_\mathcal{X}$-modules. Then $\mathcal{F}$
is quasi-coherent if and only if the following two conditions hold
\begin{enumerate}
\item $\mathcal{F}$ is locally quasi-coherent, and
\item for any morphism $\varphi : x \to y$ of $\mathcal{X}$ lying over
$f : U \to V$ the comparison map
$c_\varphi : f_{small}^*\mathcal{F}|_{V_\etale} \to
\mathcal{F}|_{U_\etale}$ of
(\ref{equation-comparison-modules}) is an isomorphism.
\end{enumerate}
\end{lemma}

\begin{proof}
Assume $\mathcal{F}$ is quasi-coherent. Then $\mathcal{F}$ is a sheaf
for the fppf topology, hence a sheaf for the \'etale topology. Moreover,
any pullback of $\mathcal{F}$ to a ringed topos is quasi-coherent, hence
the restrictions $x^*\mathcal{F}|_{U_\etale}$ are quasi-coherent.
This proves $\mathcal{F}$ is locally quasi-coherent.
Let $y$ be an object of $\mathcal{X}$ with $V = p(y)$.
We have seen that $\mathcal{X}/y = (\Sch/V)_{fppf}$. By
Descent, Proposition \ref{descent-proposition-equivalence-quasi-coherent}
it follows that $y^*\mathcal{F}$ is the quasi-coherent module
associated to a (usual) quasi-coherent module $\mathcal{F}_V$ on
the scheme $V$. Hence certainly the comparison maps
(\ref{equation-comparison-modules}) are isomorphisms.

\medskip\noindent
Conversely, suppose that $\mathcal{F}$ satisfies (1) and (2).
Let $y$ be an object of $\mathcal{X}$ with $V = p(y)$. Denote
$\mathcal{F}_V$ the quasi-coherent module on
the scheme $V$ corresponding to the restriction
$y^*\mathcal{F}|_{V_\etale}$ which is quasi-coherent by
assumption (1), see
Descent, Proposition \ref{descent-proposition-equivalence-quasi-coherent}.
Condition (2) now signifies that the restrictions
$x^*\mathcal{F}|_{U_\etale}$ for $x$ over $y$ are each
isomorphic to the (\'etale sheaf associated to the) pullback of $\mathcal{F}_V$
via the corresponding morphism of schemes $U \to V$.
Hence $y^*\mathcal{F}$ is the sheaf on $(\Sch/V)_{fppf}$
associated to $\mathcal{F}_V$. Hence it is quasi-coherent (by
Descent, Proposition \ref{descent-proposition-equivalence-quasi-coherent}
again) and we see that $\mathcal{F}$ is quasi-coherent on $\mathcal{X}$ by
Lemma \ref{lemma-characterize-quasi-coherent}.
\end{proof}

\end{document}
